% Source PDF pages 120--124; manual pages 2-91--2-95.
\section{Trajectory Calculations}\label{sec:trajectory-calculations}

The preceding sections describe the individual calculations used by TAOS. This
section shows how those pieces are organized into the complete trajectory
calculation and identifies the source-code modules that implement the workflow.
The program-module context is summarized in \cref{tab:taos-program-modules}.

\taossubhead{Main Program}

Execution begins in the \texttt{main} function in the \texttt{main} module. The
program first separates the command-line arguments into lists of table files and
problem files, as outlined in \cref{fig:main-program-flow}. A loop then processes
each problem file. For every file, TAOS creates the corresponding printout file,
although that file remains empty until the trajectory calculations for the problem
have finished.

\begin{figure}[htbp]
  \centering
  \begin{tikzpicture}[
  >=Latex,
  node distance=0.34cm,
  flowbox/.style={draw,align=center,text width=4.4cm,minimum height=0.58cm,font=\small},
  every path/.style={line width=0.65pt}
]
  \node[flowbox] (args) {Process command-line arguments};
  \node[flowbox,below=of args] (print) {Create a printout file};
  \node[flowbox,below=of print] (tables) {If first problem file, read all table files};
  \node[flowbox,below=of tables] (read) {Read problem input data};
  \node[flowbox,below=of read] (compute) {Compute trajectories};
  \node[flowbox,below=of compute] (output) {Output trajectory data};
  \node[flowbox,below=of output] (survey) {Repeat for each combination of survey values};
  \node[flowbox,below=of survey] (problem) {Repeat for each problem};
  \node[flowbox,below=of problem] (file) {Repeat for each problem file};

  \draw[->] (args)--(print);
  \draw[->] (print)--(tables);
  \draw[->] (tables)--(read);
  \draw[->] (read)--(compute);
  \draw[->] (compute)--(output);
  \draw[->] (output)--(survey);
  \draw[->] (survey)--(problem);
  \draw[->] (problem)--(file);

  \draw[->] (survey.west) -- ++(-1.0,0) |- (compute.west);
  \draw[->] (problem.west) -- ++(-1.55,0) |- (read.west);
  \draw[->] (file.west) -- ++(-2.05,0) |- (print.west);
  \draw[->] (file.south) -- ++(0,-0.45);
\end{tikzpicture}

  \caption{Main program flowchart.}
  \label{fig:main-program-flow}
\end{figure}

When the first problem file is processed, all table files are read and retained in
memory for use by that problem and every later problem file. A problem file may
contain more than one problem, so a second loop reads each problem into an
in-memory data structure.

The top level of that structure is a linked list with one entry for each vehicle or
trajectory. Each vehicle record contains its input trajectory definition, linked
segment definitions, and linked output-variable histories, as illustrated in
\cref{fig:vehicle-data-structure}. The actual program structure includes more
records and links than the figure shows; the figure identifies only the principal
relationships.

\begin{figure}[htbp]
  \centering
  \begin{tikzpicture}[
  >=Latex,
  node distance=0.42cm and 0.58cm,
  block/.style={draw,align=center,minimum height=0.68cm,font=\small},
  vehicle/.style={block,minimum width=1.65cm},
  data/.style={block,text width=1.55cm,minimum height=2.45cm},
  segment/.style={block,minimum width=1.05cm},
  output/.style={block,text width=1.4cm,minimum height=1.0cm},
  every path/.style={line width=0.65pt}
]
  \node[vehicle] (v1) {Vehicle 1};
  \node[vehicle,below=of v1] (v2) {Vehicle 2};
  \node[vehicle,below=of v2] (v3) {Vehicle 3};
  \node[below=0.17cm of v3] (vdots) {$\vdots$};
  \node[data,right=1.15cm of v2] (data) {Vehicle\\data\\structure};
  \node[segment,right=0.85cm of data.north east,anchor=north west] (s1) {Seg\\1};
  \node[segment,right=of s1] (s2) {Seg\\2};
  \node[segment,right=of s2] (s3) {Seg\\3};
  \node[right=0.18cm of s3] (sdots) {etc.};
  \node[output,below=0.55cm of s1] (o1) {Output\\variables\\at $t_1$};
  \node[output,right=of o1] (o2) {Output\\variables\\at $t_2$};
  \node[output,right=of o2] (o3) {Output\\variables\\at $t_3$};
  \node[right=0.18cm of o3] (odots) {etc.};

  \draw[->] (v1.east) -- (data.west);
  \draw[->] (v1.south) -- (v2.north);
  \draw[->] (v2.south) -- (v3.north);
  \draw[->] (v3.south) -- (vdots.north);
  \draw[->] (data.north east) -| (s1.west);
  \draw[->] (s1)--(s2);
  \draw[->] (s2)--(s3);
  \draw[->] (s3)--(sdots);
  \draw[->] (data.south east) -| (o1.west);
  \draw[->] (o1)--(o2);
  \draw[->] (o2)--(o3);
  \draw[->] (o3)--(odots);
\end{tikzpicture}

  \caption{Vehicle data structure.}
  \label{fig:vehicle-data-structure}
\end{figure}

Dynamic allocation and linked lists avoid fixed array limits on the number of
vehicles, segments, output values, or trajectory duration. In practice, the available
memory and processing time provide the limits.

After the problem data are read, TAOS computes all trajectories and then writes the
trajectory data to the printout and other requested output files. Because trajectory
output is not emitted until all trajectories have been calculated, a premature program
termination can leave the output files empty. When a survey is defined, the complete
trajectory calculation is repeated for every combination of survey values.

\taossubhead{Trajectory Calculation Loop}

The \texttt{compute\_trajectories} function in the \texttt{path} module expands the
``compute trajectories'' step of \cref{fig:main-program-flow}. Its logic is shown in
\cref{fig:trajectory-calculation-flow}.

\begin{figure}[htbp]
  \centering
  \begin{tikzpicture}[
  >=Latex,
  node distance=0.45cm,
  flowbox/.style={draw,align=center,text width=5.2cm,minimum height=0.72cm,font=\small},
  decision/.style={draw,rounded corners=8pt,align=center,text width=4.2cm,minimum height=0.72cm,font=\small},
  every path/.style={line width=0.7pt}
]
  \node[flowbox] (init) {Initialize starting trajectories};
  \node[flowbox,below=of init] (dt) {Compute a $\Delta t$};
  \node[flowbox,below=of dt] (integrate) {Integrate all trajectories one time step};
  \node[flowbox,below=of integrate] (final) {If a segment final condition is hit, change segment or end trajectory};
  \node[flowbox,below=of final] (search) {If a search function is hit, restart trajectories as necessary for the search};
  \node[flowbox,below=of search] (new) {Check for starting new trajectories};
  \node[decision,below=of new] (done) {All trajectories completed?};

  \draw[->] (init)--(dt);
  \draw[->] (dt)--(integrate);
  \draw[->] (integrate)--(final);
  \draw[->] (final)--(search);
  \draw[->] (search)--(new);
  \draw[->] (new)--(done);
  \draw[->] (done.west) -- ++(-1.1,0) |- (dt.west)
    node[pos=0.13,above,font=\scriptsize] {no};
  \draw[->] (done.south) -- ++(0,-0.48)
    node[below,font=\scriptsize] {yes};
\end{tikzpicture}

  \caption{Trajectory calculation flowchart.}
  \label{fig:trajectory-calculation-flow}
\end{figure}

The calculation begins by locating the starting vehicle, defined as the vehicle with
the minimum initial time. A trajectory may be initialized from directly specified
initial conditions or from a state on an existing trajectory, but at least one trajectory
must have direct initial conditions so that a starting trajectory exists.

TAOS maintains an active/inactive status for every trajectory. A trajectory is active
when it has been initialized and is currently being integrated. It is inactive before
initialization and after completion.

After the starting trajectories are initialized, TAOS chooses the smallest integration
step requested by any active vehicle. The step is shortened when necessary to land
exactly on a print time or on a time value appearing in a guidance-rule table. All
active trajectories are synchronized and advanced with the same step
$\Delta t$. The step selection is implemented by
\texttt{get\_next\_time\_step} in the \texttt{path} module.

Each active vehicle is integrated separately. The vehicles are not assembled into one
large system of differential equations. The fourth-order Runge--Kutta procedure in
\texttt{integration\_step} advances one vehicle through the selected time step, using
the method described in \cref{sec:numerical-integration}.

\taossubhead{Segment Final Conditions}

After every integration step, \texttt{check\_final\_conditions} in the \texttt{path}
module tests the final conditions for each active segment. If a final-condition
boundary was crossed, TAOS repeats the last step while varying $\Delta t$ until the
condition is met within a small tolerance. The resulting one-dimensional search uses
the secant method described in \cref{sec:search-methods}.

More than one segment can appear to end within the same integration step. If the
final conditions occur at exactly the same time, the segments end simultaneously.
If they occur at different times, the step is repeatedly halved until only one final
condition is crossed; TAOS then applies the secant search to that condition.

\taossubhead{Search and Optimization Loops}

Search and optimization loops vary one or more problem inputs to meet a trajectory
goal. The goal is normally a simple expression, often an output variable, evaluated
at the end of a segment. TAOS permits search and optimization objective functions
only at segment boundaries.

At each segment end, \texttt{search\_function\_hit} in the \texttt{search} module
determines whether a search or optimization function is associated with that point.
When one is found, TAOS modifies the designated input parameters according to the
selected algorithm and restarts the affected trajectories. The one-dimensional search
methods are described in \cref{sec:search-methods}; the optimization method is
described in \cref{sec:optimization}.

A restart begins immediately before the earliest modified input parameter. Portions
of the trajectories that precede every modified value remain valid and are not
recomputed. This reduces the calculation required by nested searches and
optimization loops.

\taossubhead{Multiple Trajectories}

When a segment ends, its trajectory either continues with another segment or is
marked complete. During segment initialization, \texttt{add\_new\_vehicles} in the
\texttt{path} module examines inactive trajectories. Any trajectory whose initial
state depends on the segment that has just been reached is initialized and added to
the active list. The calculation loop ends only after every trajectory has been marked
complete.
